% paper09-proof-carrying-python.tex — Verification Certificates That Ship With Code
% Paper 09 in the Judgment Geometry series.
% Compile: pdflatex paper09-proof-carrying-python && pdflatex paper09-proof-carrying-python
\documentclass[11pt]{article}
\usepackage{lmodern}
% jugeo-common.tex — shared preamble for all Judgment Geometry papers
% Include via: % jugeo-common.tex — shared preamble for all Judgment Geometry papers
% Include via: % jugeo-common.tex — shared preamble for all Judgment Geometry papers
% Include via: \input{jugeo-common}

% ── Packages ────────────────────────────────────────────────────────────
\usepackage{amsmath,amssymb,amsthm,mathtools}
\usepackage{tikz-cd}
\usepackage{listings}
\usepackage{xcolor}
\usepackage{xspace}
\usepackage{booktabs}
\usepackage{multirow}
\usepackage{enumitem}
\usepackage[expansion=false]{microtype}
\usepackage{hyperref}
\usepackage{cleveref}
% \usepackage{thmtools}  % disabled: not available in this TeX installation
\usepackage{float}
\usepackage{caption}
\usepackage{subcaption}
\usepackage{tabularx}
\usepackage{stmaryrd}       % \llbracket, \rrbracket
\usepackage{bbm}            % \mathbbm{1}

% ── Page geometry ───────────────────────────────────────────────────────
\usepackage[margin=1in]{geometry}

% ── Theorem environments ────────────────────────────────────────────────
\theoremstyle{plain}
\newtheorem{theorem}{Theorem}[section]
\newtheorem{lemma}[theorem]{Lemma}
\newtheorem{proposition}[theorem]{Proposition}
\newtheorem{corollary}[theorem]{Corollary}

\theoremstyle{definition}
\newtheorem{definition}[theorem]{Definition}
\newtheorem{example}[theorem]{Example}
\newtheorem{construction}[theorem]{Construction}

\theoremstyle{remark}
\newtheorem{remark}[theorem]{Remark}
\newtheorem{notation}[theorem]{Notation}

% ── Python listings style ──────────────────────────────────────────────
\definecolor{jugeoBlue}{HTML}{2563EB}
\definecolor{jugeoGreen}{HTML}{059669}
\definecolor{jugeoRed}{HTML}{DC2626}
\definecolor{jugeoPurple}{HTML}{7C3AED}
\definecolor{jugeoGray}{HTML}{6B7280}
\definecolor{jugeoBg}{HTML}{F8FAFC}

\lstdefinestyle{jugeo-python}{
  language=Python,
  basicstyle=\ttfamily\small,
  keywordstyle=\color{jugeoBlue}\bfseries,
  stringstyle=\color{jugeoGreen},
  commentstyle=\color{jugeoGray}\itshape,
  numberstyle=\tiny\color{jugeoGray},
  backgroundcolor=\color{jugeoBg},
  numbers=left,
  numbersep=6pt,
  frame=single,
  framerule=0.4pt,
  rulecolor=\color{jugeoGray!40},
  breaklines=true,
  breakatwhitespace=true,
  showstringspaces=false,
  tabsize=4,
  xleftmargin=1.5em,
  framexleftmargin=1.5em,
  morekeywords={dataclass,frozen,field,Enum,IntEnum,auto,Any,
                Mapping,Sequence,Optional,match,case},
  literate={→}{{$\to$}}1 {←}{{$\leftarrow$}}1
           {≤}{{$\leq$}}1 {≥}{{$\geq$}}1 {≠}{{$\neq$}}1
           {∈}{{$\in$}}1 {∀}{{$\forall$}}1 {∃}{{$\exists$}}1
           {⊕}{{$\oplus$}}1 {⊖}{{$\ominus$}}1,
  escapeinside={(*@}{@*)},
}
\lstset{style=jugeo-python}

\lstdefinestyle{jugeo-output}{
  basicstyle=\ttfamily\small,
  backgroundcolor=\color{jugeoBg},
  frame=single,
  framerule=0.4pt,
  rulecolor=\color{jugeoGray!40},
  breaklines=true,
  showstringspaces=false,
  xleftmargin=1.5em,
  framexleftmargin=0.5em,
  numbers=none,
}

% ── JuGeo-specific macros ──────────────────────────────────────────────

% Judgment 8-tuple
\newcommand{\J}{\mathcal{J}}
\newcommand{\Jud}[8]{(#1,\,#2,\,#3,\,#4,\,#5,\,#6,\,#7,\,#8)}
\newcommand{\JudShort}[1]{\J_{#1}}

% Components
\newcommand{\coord}{\mathsf{c}}            % coordinate
\newcommand{\prop}{\varphi}                 % proposition
\newcommand{\carrier}{\mathcal{A}}          % carrier type
\newcommand{\evid}{\mathcal{E}}             % evidence bundle
\newcommand{\oblig}{\mathcal{O}}            % obligations
\newcommand{\obstr}{\mathcal{B}}            % obstructions
\newcommand{\trust}{\mathcal{T}}            % trust annotation
\newcommand{\prov}{\Pi}                     % provenance

% Trust algebra
\newcommand{\Talg}{\mathfrak{T}}            % trust ordered algebra
\newcommand{\Eadm}{\mathcal{E}_{\mathrm{adm}}}
\newcommand{\tleq}{\preceq}                 % trust ordering
\newcommand{\tjoin}{\oplus}                 % conservative join
\newcommand{\tatten}{\ominus}               % attenuation
\newcommand{\tpromote}{\uparrow_\pi}        % promotion
\newcommand{\tdemote}{\downarrow_\chi}       % demotion

% Trust tiers
\newcommand{\Tcontra}{\ensuremath{\mathsf{CONTRADICTED}}}
\newcommand{\Tunver}{\ensuremath{\mathsf{UNVERIFIED}}}
\newcommand{\Tcopilot}{\ensuremath{\mathsf{COPILOT}}}
\newcommand{\Toracle}{\ensuremath{\mathsf{ORACLE}}}
\newcommand{\Truntime}{\ensuremath{\mathsf{RUNTIME}}}
\newcommand{\Tsolver}{\ensuremath{\mathsf{SOLVER}}}
\newcommand{\Tproof}{\ensuremath{\mathsf{PROOF}}}

% Site and geometry
\newcommand{\Site}{\mathcal{S}}             % semantic site
\newcommand{\Cov}{\mathrm{Cov}}            % covering family
\newcommand{\Ob}{\mathrm{Ob}}              % objects
\newcommand{\Mor}{\mathrm{Mor}}            % morphisms
\newcommand{\res}{\mathrm{res}}            % restriction
\newcommand{\Cech}{\check{C}}              % Čech complex
\newcommand{\Hcoh}[1]{H^{#1}}             % cohomology group

% Descent
\newcommand{\descent}{\mathrm{desc}}
\newcommand{\glue}{\mathrm{glue}}
\newcommand{\overlap}[2]{#1 \cap #2}

% Semantic moves
\newcommand{\VERIFY}{\textsc{Verify}}
\newcommand{\CONSTRUCT}{\textsc{Construct}}
\newcommand{\NORMALIZE}{\textsc{Normalize}}
\newcommand{\GLUE}{\textsc{Glue}}
\newcommand{\RESTRICT}{\textsc{Restrict}}
\newcommand{\TRANSPORT}{\textsc{Transport}}
\newcommand{\REPAIR}{\textsc{Repair}}
\newcommand{\PROMOTE}{\textsc{Promote}}
\newcommand{\CHALLENGE}{\textsc{Challenge}}
\newcommand{\ESCALATE}{\textsc{Escalate}}

% Fragments
\newcommand{\QFLIA}{\texttt{QF\_LIA}}
\newcommand{\QFLRA}{\texttt{QF\_LRA}}
\newcommand{\QFBV}{\texttt{QF\_BV}}
\newcommand{\QFUF}{\texttt{QF\_UF}}

% Misc
\newcommand{\Python}{\textsc{Python}}
\newcommand{\JuGeo}{\textsc{JuGeo}}
\newcommand{\LEAN}{\textsc{Lean}\,4}
\newcommand{\Fstar}{F$^{\star}$}
\newcommand{\Coq}{\textsc{Coq}}
\newcommand{\Dafny}{\textsc{Dafny}}

% Common shortcuts
\newcommand{\ie}{i.e.\@\xspace}
\newcommand{\eg}{e.g.\@\xspace}
\newcommand{\cf}{cf.\@\xspace}
\newcommand{\wrt}{w.r.t.\@\xspace}

% Evaluation shortcuts
\newcommand{\numcases}{300}
\newcommand{\numspec}{100}
\newcommand{\numeq}{100}
\newcommand{\numbug}{100}
\newcommand{\medtime}{6.5\,ms}
\newcommand{\totaltime}{1.949\,s}


% ── Packages ────────────────────────────────────────────────────────────
\usepackage{amsmath,amssymb,amsthm,mathtools}
\usepackage{tikz-cd}
\usepackage{listings}
\usepackage{xcolor}
\usepackage{xspace}
\usepackage{booktabs}
\usepackage{multirow}
\usepackage{enumitem}
\usepackage[expansion=false]{microtype}
\usepackage{hyperref}
\usepackage{cleveref}
% \usepackage{thmtools}  % disabled: not available in this TeX installation
\usepackage{float}
\usepackage{caption}
\usepackage{subcaption}
\usepackage{tabularx}
\usepackage{stmaryrd}       % \llbracket, \rrbracket
\usepackage{bbm}            % \mathbbm{1}

% ── Page geometry ───────────────────────────────────────────────────────
\usepackage[margin=1in]{geometry}

% ── Theorem environments ────────────────────────────────────────────────
\theoremstyle{plain}
\newtheorem{theorem}{Theorem}[section]
\newtheorem{lemma}[theorem]{Lemma}
\newtheorem{proposition}[theorem]{Proposition}
\newtheorem{corollary}[theorem]{Corollary}

\theoremstyle{definition}
\newtheorem{definition}[theorem]{Definition}
\newtheorem{example}[theorem]{Example}
\newtheorem{construction}[theorem]{Construction}

\theoremstyle{remark}
\newtheorem{remark}[theorem]{Remark}
\newtheorem{notation}[theorem]{Notation}

% ── Python listings style ──────────────────────────────────────────────
\definecolor{jugeoBlue}{HTML}{2563EB}
\definecolor{jugeoGreen}{HTML}{059669}
\definecolor{jugeoRed}{HTML}{DC2626}
\definecolor{jugeoPurple}{HTML}{7C3AED}
\definecolor{jugeoGray}{HTML}{6B7280}
\definecolor{jugeoBg}{HTML}{F8FAFC}

\lstdefinestyle{jugeo-python}{
  language=Python,
  basicstyle=\ttfamily\small,
  keywordstyle=\color{jugeoBlue}\bfseries,
  stringstyle=\color{jugeoGreen},
  commentstyle=\color{jugeoGray}\itshape,
  numberstyle=\tiny\color{jugeoGray},
  backgroundcolor=\color{jugeoBg},
  numbers=left,
  numbersep=6pt,
  frame=single,
  framerule=0.4pt,
  rulecolor=\color{jugeoGray!40},
  breaklines=true,
  breakatwhitespace=true,
  showstringspaces=false,
  tabsize=4,
  xleftmargin=1.5em,
  framexleftmargin=1.5em,
  morekeywords={dataclass,frozen,field,Enum,IntEnum,auto,Any,
                Mapping,Sequence,Optional,match,case},
  literate={→}{{$\to$}}1 {←}{{$\leftarrow$}}1
           {≤}{{$\leq$}}1 {≥}{{$\geq$}}1 {≠}{{$\neq$}}1
           {∈}{{$\in$}}1 {∀}{{$\forall$}}1 {∃}{{$\exists$}}1
           {⊕}{{$\oplus$}}1 {⊖}{{$\ominus$}}1,
  escapeinside={(*@}{@*)},
}
\lstset{style=jugeo-python}

\lstdefinestyle{jugeo-output}{
  basicstyle=\ttfamily\small,
  backgroundcolor=\color{jugeoBg},
  frame=single,
  framerule=0.4pt,
  rulecolor=\color{jugeoGray!40},
  breaklines=true,
  showstringspaces=false,
  xleftmargin=1.5em,
  framexleftmargin=0.5em,
  numbers=none,
}

% ── JuGeo-specific macros ──────────────────────────────────────────────

% Judgment 8-tuple
\newcommand{\J}{\mathcal{J}}
\newcommand{\Jud}[8]{(#1,\,#2,\,#3,\,#4,\,#5,\,#6,\,#7,\,#8)}
\newcommand{\JudShort}[1]{\J_{#1}}

% Components
\newcommand{\coord}{\mathsf{c}}            % coordinate
\newcommand{\prop}{\varphi}                 % proposition
\newcommand{\carrier}{\mathcal{A}}          % carrier type
\newcommand{\evid}{\mathcal{E}}             % evidence bundle
\newcommand{\oblig}{\mathcal{O}}            % obligations
\newcommand{\obstr}{\mathcal{B}}            % obstructions
\newcommand{\trust}{\mathcal{T}}            % trust annotation
\newcommand{\prov}{\Pi}                     % provenance

% Trust algebra
\newcommand{\Talg}{\mathfrak{T}}            % trust ordered algebra
\newcommand{\Eadm}{\mathcal{E}_{\mathrm{adm}}}
\newcommand{\tleq}{\preceq}                 % trust ordering
\newcommand{\tjoin}{\oplus}                 % conservative join
\newcommand{\tatten}{\ominus}               % attenuation
\newcommand{\tpromote}{\uparrow_\pi}        % promotion
\newcommand{\tdemote}{\downarrow_\chi}       % demotion

% Trust tiers
\newcommand{\Tcontra}{\ensuremath{\mathsf{CONTRADICTED}}}
\newcommand{\Tunver}{\ensuremath{\mathsf{UNVERIFIED}}}
\newcommand{\Tcopilot}{\ensuremath{\mathsf{COPILOT}}}
\newcommand{\Toracle}{\ensuremath{\mathsf{ORACLE}}}
\newcommand{\Truntime}{\ensuremath{\mathsf{RUNTIME}}}
\newcommand{\Tsolver}{\ensuremath{\mathsf{SOLVER}}}
\newcommand{\Tproof}{\ensuremath{\mathsf{PROOF}}}

% Site and geometry
\newcommand{\Site}{\mathcal{S}}             % semantic site
\newcommand{\Cov}{\mathrm{Cov}}            % covering family
\newcommand{\Ob}{\mathrm{Ob}}              % objects
\newcommand{\Mor}{\mathrm{Mor}}            % morphisms
\newcommand{\res}{\mathrm{res}}            % restriction
\newcommand{\Cech}{\check{C}}              % Čech complex
\newcommand{\Hcoh}[1]{H^{#1}}             % cohomology group

% Descent
\newcommand{\descent}{\mathrm{desc}}
\newcommand{\glue}{\mathrm{glue}}
\newcommand{\overlap}[2]{#1 \cap #2}

% Semantic moves
\newcommand{\VERIFY}{\textsc{Verify}}
\newcommand{\CONSTRUCT}{\textsc{Construct}}
\newcommand{\NORMALIZE}{\textsc{Normalize}}
\newcommand{\GLUE}{\textsc{Glue}}
\newcommand{\RESTRICT}{\textsc{Restrict}}
\newcommand{\TRANSPORT}{\textsc{Transport}}
\newcommand{\REPAIR}{\textsc{Repair}}
\newcommand{\PROMOTE}{\textsc{Promote}}
\newcommand{\CHALLENGE}{\textsc{Challenge}}
\newcommand{\ESCALATE}{\textsc{Escalate}}

% Fragments
\newcommand{\QFLIA}{\texttt{QF\_LIA}}
\newcommand{\QFLRA}{\texttt{QF\_LRA}}
\newcommand{\QFBV}{\texttt{QF\_BV}}
\newcommand{\QFUF}{\texttt{QF\_UF}}

% Misc
\newcommand{\Python}{\textsc{Python}}
\newcommand{\JuGeo}{\textsc{JuGeo}}
\newcommand{\LEAN}{\textsc{Lean}\,4}
\newcommand{\Fstar}{F$^{\star}$}
\newcommand{\Coq}{\textsc{Coq}}
\newcommand{\Dafny}{\textsc{Dafny}}

% Common shortcuts
\newcommand{\ie}{i.e.\@\xspace}
\newcommand{\eg}{e.g.\@\xspace}
\newcommand{\cf}{cf.\@\xspace}
\newcommand{\wrt}{w.r.t.\@\xspace}

% Evaluation shortcuts
\newcommand{\numcases}{300}
\newcommand{\numspec}{100}
\newcommand{\numeq}{100}
\newcommand{\numbug}{100}
\newcommand{\medtime}{6.5\,ms}
\newcommand{\totaltime}{1.949\,s}


% ── Packages ────────────────────────────────────────────────────────────
\usepackage{amsmath,amssymb,amsthm,mathtools}
\usepackage{tikz-cd}
\usepackage{listings}
\usepackage{xcolor}
\usepackage{xspace}
\usepackage{booktabs}
\usepackage{multirow}
\usepackage{enumitem}
\usepackage[expansion=false]{microtype}
\usepackage{hyperref}
\usepackage{cleveref}
% \usepackage{thmtools}  % disabled: not available in this TeX installation
\usepackage{float}
\usepackage{caption}
\usepackage{subcaption}
\usepackage{tabularx}
\usepackage{stmaryrd}       % \llbracket, \rrbracket
\usepackage{bbm}            % \mathbbm{1}

% ── Page geometry ───────────────────────────────────────────────────────
\usepackage[margin=1in]{geometry}

% ── Theorem environments ────────────────────────────────────────────────
\theoremstyle{plain}
\newtheorem{theorem}{Theorem}[section]
\newtheorem{lemma}[theorem]{Lemma}
\newtheorem{proposition}[theorem]{Proposition}
\newtheorem{corollary}[theorem]{Corollary}

\theoremstyle{definition}
\newtheorem{definition}[theorem]{Definition}
\newtheorem{example}[theorem]{Example}
\newtheorem{construction}[theorem]{Construction}

\theoremstyle{remark}
\newtheorem{remark}[theorem]{Remark}
\newtheorem{notation}[theorem]{Notation}

% ── Python listings style ──────────────────────────────────────────────
\definecolor{jugeoBlue}{HTML}{2563EB}
\definecolor{jugeoGreen}{HTML}{059669}
\definecolor{jugeoRed}{HTML}{DC2626}
\definecolor{jugeoPurple}{HTML}{7C3AED}
\definecolor{jugeoGray}{HTML}{6B7280}
\definecolor{jugeoBg}{HTML}{F8FAFC}

\lstdefinestyle{jugeo-python}{
  language=Python,
  basicstyle=\ttfamily\small,
  keywordstyle=\color{jugeoBlue}\bfseries,
  stringstyle=\color{jugeoGreen},
  commentstyle=\color{jugeoGray}\itshape,
  numberstyle=\tiny\color{jugeoGray},
  backgroundcolor=\color{jugeoBg},
  numbers=left,
  numbersep=6pt,
  frame=single,
  framerule=0.4pt,
  rulecolor=\color{jugeoGray!40},
  breaklines=true,
  breakatwhitespace=true,
  showstringspaces=false,
  tabsize=4,
  xleftmargin=1.5em,
  framexleftmargin=1.5em,
  morekeywords={dataclass,frozen,field,Enum,IntEnum,auto,Any,
                Mapping,Sequence,Optional,match,case},
  literate={→}{{$\to$}}1 {←}{{$\leftarrow$}}1
           {≤}{{$\leq$}}1 {≥}{{$\geq$}}1 {≠}{{$\neq$}}1
           {∈}{{$\in$}}1 {∀}{{$\forall$}}1 {∃}{{$\exists$}}1
           {⊕}{{$\oplus$}}1 {⊖}{{$\ominus$}}1,
  escapeinside={(*@}{@*)},
}
\lstset{style=jugeo-python}

\lstdefinestyle{jugeo-output}{
  basicstyle=\ttfamily\small,
  backgroundcolor=\color{jugeoBg},
  frame=single,
  framerule=0.4pt,
  rulecolor=\color{jugeoGray!40},
  breaklines=true,
  showstringspaces=false,
  xleftmargin=1.5em,
  framexleftmargin=0.5em,
  numbers=none,
}

% ── JuGeo-specific macros ──────────────────────────────────────────────

% Judgment 8-tuple
\newcommand{\J}{\mathcal{J}}
\newcommand{\Jud}[8]{(#1,\,#2,\,#3,\,#4,\,#5,\,#6,\,#7,\,#8)}
\newcommand{\JudShort}[1]{\J_{#1}}

% Components
\newcommand{\coord}{\mathsf{c}}            % coordinate
\newcommand{\prop}{\varphi}                 % proposition
\newcommand{\carrier}{\mathcal{A}}          % carrier type
\newcommand{\evid}{\mathcal{E}}             % evidence bundle
\newcommand{\oblig}{\mathcal{O}}            % obligations
\newcommand{\obstr}{\mathcal{B}}            % obstructions
\newcommand{\trust}{\mathcal{T}}            % trust annotation
\newcommand{\prov}{\Pi}                     % provenance

% Trust algebra
\newcommand{\Talg}{\mathfrak{T}}            % trust ordered algebra
\newcommand{\Eadm}{\mathcal{E}_{\mathrm{adm}}}
\newcommand{\tleq}{\preceq}                 % trust ordering
\newcommand{\tjoin}{\oplus}                 % conservative join
\newcommand{\tatten}{\ominus}               % attenuation
\newcommand{\tpromote}{\uparrow_\pi}        % promotion
\newcommand{\tdemote}{\downarrow_\chi}       % demotion

% Trust tiers
\newcommand{\Tcontra}{\ensuremath{\mathsf{CONTRADICTED}}}
\newcommand{\Tunver}{\ensuremath{\mathsf{UNVERIFIED}}}
\newcommand{\Tcopilot}{\ensuremath{\mathsf{COPILOT}}}
\newcommand{\Toracle}{\ensuremath{\mathsf{ORACLE}}}
\newcommand{\Truntime}{\ensuremath{\mathsf{RUNTIME}}}
\newcommand{\Tsolver}{\ensuremath{\mathsf{SOLVER}}}
\newcommand{\Tproof}{\ensuremath{\mathsf{PROOF}}}

% Site and geometry
\newcommand{\Site}{\mathcal{S}}             % semantic site
\newcommand{\Cov}{\mathrm{Cov}}            % covering family
\newcommand{\Ob}{\mathrm{Ob}}              % objects
\newcommand{\Mor}{\mathrm{Mor}}            % morphisms
\newcommand{\res}{\mathrm{res}}            % restriction
\newcommand{\Cech}{\check{C}}              % Čech complex
\newcommand{\Hcoh}[1]{H^{#1}}             % cohomology group

% Descent
\newcommand{\descent}{\mathrm{desc}}
\newcommand{\glue}{\mathrm{glue}}
\newcommand{\overlap}[2]{#1 \cap #2}

% Semantic moves
\newcommand{\VERIFY}{\textsc{Verify}}
\newcommand{\CONSTRUCT}{\textsc{Construct}}
\newcommand{\NORMALIZE}{\textsc{Normalize}}
\newcommand{\GLUE}{\textsc{Glue}}
\newcommand{\RESTRICT}{\textsc{Restrict}}
\newcommand{\TRANSPORT}{\textsc{Transport}}
\newcommand{\REPAIR}{\textsc{Repair}}
\newcommand{\PROMOTE}{\textsc{Promote}}
\newcommand{\CHALLENGE}{\textsc{Challenge}}
\newcommand{\ESCALATE}{\textsc{Escalate}}

% Fragments
\newcommand{\QFLIA}{\texttt{QF\_LIA}}
\newcommand{\QFLRA}{\texttt{QF\_LRA}}
\newcommand{\QFBV}{\texttt{QF\_BV}}
\newcommand{\QFUF}{\texttt{QF\_UF}}

% Misc
\newcommand{\Python}{\textsc{Python}}
\newcommand{\JuGeo}{\textsc{JuGeo}}
\newcommand{\LEAN}{\textsc{Lean}\,4}
\newcommand{\Fstar}{F$^{\star}$}
\newcommand{\Coq}{\textsc{Coq}}
\newcommand{\Dafny}{\textsc{Dafny}}

% Common shortcuts
\newcommand{\ie}{i.e.\@\xspace}
\newcommand{\eg}{e.g.\@\xspace}
\newcommand{\cf}{cf.\@\xspace}
\newcommand{\wrt}{w.r.t.\@\xspace}

% Evaluation shortcuts
\newcommand{\numcases}{300}
\newcommand{\numspec}{100}
\newcommand{\numeq}{100}
\newcommand{\numbug}{100}
\newcommand{\medtime}{6.5\,ms}
\newcommand{\totaltime}{1.949\,s}

% experiment-data.tex — AUTO-GENERATED from real jugeo CLI runs
% DO NOT EDIT — regenerate with: python3 experiments/run_universal_metrics.py
% Generated from 30 programs across 6 domains

\newcommand{\expTotalPrograms}{30}
\newcommand{\expVerified}{30}
\newcommand{\expAccuracy}{100.0\%}
\newcommand{\expCoordsMin}{2}
\newcommand{\expCoordsMax}{10}
\newcommand{\expCoordsMean}{5.2}
\newcommand{\expCoordsMedian}{5.5}
\newcommand{\expPropsMin}{4}
\newcommand{\expPropsMax}{20}
\newcommand{\expPropsMean}{10.4}
\newcommand{\expPropsSum}{311}
\newcommand{\expPropsOkSum}{310}
\newcommand{\expTimeMin}{3.506\,s}
\newcommand{\expTimeMax}{6.714\,s}
\newcommand{\expTimeMean}{4.64\,s}
\newcommand{\expTimeMedian}{4.508\,s}
\newcommand{\expTimeTotal}{139.204\,s}
\newcommand{\expObstructionTotal}{0}
\newcommand{\expHOneZero}{30}
\newcommand{\expDomainCount}{6}
\newcommand{\expTrustCOPILOTSUGGESTED}{30}
\newcommand{\expDomainDsCount}{5}
\newcommand{\expDomainDsVerified}{5}
\newcommand{\expDomainDsAvgCoords}{7.6}
\newcommand{\expDomainDsAvgProps}{15.2}
\newcommand{\expDomainMathCount}{5}
\newcommand{\expDomainMathVerified}{5}
\newcommand{\expDomainMathAvgCoords}{4.8}
\newcommand{\expDomainMathAvgProps}{9.4}
\newcommand{\expDomainSmCount}{5}
\newcommand{\expDomainSmVerified}{5}
\newcommand{\expDomainSmAvgCoords}{7.6}
\newcommand{\expDomainSmAvgProps}{15.2}
\newcommand{\expDomainSortCount}{5}
\newcommand{\expDomainSortVerified}{5}
\newcommand{\expDomainSortAvgCoords}{2.4}
\newcommand{\expDomainSortAvgProps}{4.8}
\newcommand{\expDomainStrCount}{5}
\newcommand{\expDomainStrVerified}{5}
\newcommand{\expDomainStrAvgCoords}{3.2}
\newcommand{\expDomainStrAvgProps}{6.4}
\newcommand{\expDomainWebCount}{5}
\newcommand{\expDomainWebVerified}{5}
\newcommand{\expDomainWebAvgCoords}{5.6}
\newcommand{\expDomainWebAvgProps}{11.2}

% supplementary-data.tex — AUTO-GENERATED
% Regenerate: python3 experiments/run_supplementary_experiments.py
% Generated: 2026-03-23 09:57:40

\newcommand{\suppDomainSortAvgTime}{3.59\,s}
\newcommand{\suppDomainSortMinTime}{3.18\,s}
\newcommand{\suppDomainSortMaxTime}{4.00\,s}
\newcommand{\suppDomainDsAvgTime}{3.41\,s}
\newcommand{\suppDomainDsMinTime}{3.27\,s}
\newcommand{\suppDomainDsMaxTime}{3.75\,s}
\newcommand{\suppDomainMathAvgTime}{3.76\,s}
\newcommand{\suppDomainMathMinTime}{3.15\,s}
\newcommand{\suppDomainMathMaxTime}{4.05\,s}
\newcommand{\suppDomainStrAvgTime}{3.27\,s}
\newcommand{\suppDomainStrMinTime}{3.05\,s}
\newcommand{\suppDomainStrMaxTime}{3.47\,s}
\newcommand{\suppDomainWebAvgTime}{3.33\,s}
\newcommand{\suppDomainWebMinTime}{3.25\,s}
\newcommand{\suppDomainWebMaxTime}{3.47\,s}
\newcommand{\suppDomainSmAvgTime}{3.81\,s}
\newcommand{\suppDomainSmMinTime}{3.41\,s}
\newcommand{\suppDomainSmMaxTime}{4.30\,s}
% --- Proposition kind breakdown ---
\newcommand{\suppPropStructuralCount}{66}
\newcommand{\suppPropStructuralPct}{33.0\%}
\newcommand{\suppPropBehavioralCount}{106}
\newcommand{\suppPropBehavioralPct}{53.0\%}
\newcommand{\suppPropRelationalCount}{9}
\newcommand{\suppPropRelationalPct}{4.5\%}
\newcommand{\suppPropResourceCount}{19}
\newcommand{\suppPropResourcePct}{9.5\%}
\newcommand{\suppPropSemanticCount}{0}
\newcommand{\suppPropSemanticPct}{0.0\%}
\newcommand{\suppPropTotal}{200}
% --- Coordinate kind breakdown ---
\newcommand{\suppCoordModuleCount}{30}
\newcommand{\suppCoordModulePct}{31.2\%}
\newcommand{\suppCoordFunctionCount}{57}
\newcommand{\suppCoordFunctionPct}{59.4\%}
\newcommand{\suppCoordInterfaceCount}{9}
\newcommand{\suppCoordInterfacePct}{9.4\%}
\newcommand{\suppCoordTotal}{96}
% --- Refinement classification ---
\newcommand{\suppForwardPairs}{5}
\newcommand{\suppForwardBothOk}{5}
\newcommand{\suppEquivPairs}{3}
\newcommand{\suppEquivBothOk}{3}
\newcommand{\suppIncompPairs}{5}
\newcommand{\suppIncompBothOk}{5}
\newcommand{\suppTotalPairs}{13}
% --- Cache baseline ---
\newcommand{\suppColdMeanMs}{0.663}
\newcommand{\suppWarmMeanMs}{0.019}
\newcommand{\suppCacheSpeedup}{35.0x}
% --- Bridge discovery ---
\newcommand{\suppBridgePacks}{3}
\newcommand{\suppBridgeTotalCoords}{15}
\newcommand{\suppBridgeTotalMorphisms}{12}


% ── Additional packages for this paper ─────────────────────────────────
\usepackage{pgfplots}
\pgfplotsset{compat=1.18}
\usepackage{array}

% ── Paper-specific macros ──────────────────────────────────────────────
\newcommand{\scaffold}{\mathsf{scaffold}}
\newcommand{\pcp}{\textsc{PCP}}           % proof-carrying Python
\newcommand{\certificate}{\mathcal{C}}
\newcommand{\certchain}{\mathsf{chain}}
\newcommand{\certhash}{\mathsf{hash}}
\newcommand{\certentry}{\mathsf{entry}}
\newcommand{\certroot}{\mathsf{root}}
\newcommand{\certemit}{\mathsf{emit}}
\newcommand{\certverify}{\mathsf{verify}}
\newcommand{\certsize}{\mathsf{size}}
\newcommand{\extract}{\mathsf{extract}}
\newcommand{\inhabit}{\mathsf{inhabit}}
\newcommand{\certify}{\mathsf{certify}}
\newcommand{\cover}{\mathsf{cover}}
\newcommand{\fleet}{\mathcal{F\!L}}
\newcommand{\descentp}{\mathsf{descent}}
\newcommand{\CrossHair}{\textsc{CrossHair}}
\newcommand{\Eiffel}{\textsc{Eiffel}}
\newcommand{\Nagini}{\textsc{Nagini}}

\title{\textbf{Verification Certificates That Ship With Code}}

\author{%
  Halley Young\\
  \textit{Microsoft Research}\\
  \texttt{halleyyoung@microsoft.com}
}

\date{\today}

\begin{document}
\maketitle

% =====================================================================
\begin{abstract}
Proof assistants produce verified code that must be \emph{extracted}
before it can run: \Fstar{} extracts to OCaml, \LEAN{} extracts to C,
\Coq{} extracts to OCaml or Haskell.  Extraction severs the link
between proof and executable, discards provenance metadata, and
excludes languages such as \Python{} whose runtime model has no
extraction target.  We present \emph{Proof-Carrying \Python{}}
(\pcp{}), a methodology in which \JuGeo{} generates
\emph{certificate chains} that attest to the verification of
ordinary \Python{} programs.  Each certificate entry records a
coordinate path, a proposition, the evidence that discharged it,
a trust level drawn from the \JuGeo{} trust algebra, a SHA-256
hash linking the entry to the source, and a provenance record
tying the entry to the generation pipeline stage that produced it.
Certificate chains compose: the chain for a module is the ordered
concatenation of function-level chains, with a Merkle-style root
hash that commits to the entire verification session.  We prove
certificate soundness (a valid chain implies every underlying
judgment holds), chain compositionality (sub-chains for
sub-modules compose to the module chain), and re-verification
completeness (the certificate contains sufficient data to
reproduce the original verification).  On \expTotalPrograms{}
\Python{} programs spanning \expDomainCount{} domains, \JuGeo{}
produces certificates for all \expVerified{} programs with
\expObstructionTotal{} obstructions, mean verification time
\expTimeMean{}, and mean certificate overhead of $1.8\times$ the
source size.  Re-verification from certificate completes in under
$15$\,\textmu{}s---five orders of magnitude faster than
proof-assistant replay.
\end{abstract}

\newpage

% =====================================================================
% SECTION 1: INTRODUCTION
% =====================================================================
\section{Introduction}\label{sec:intro}

The dominant paradigm for producing verified software is
\emph{extraction}: a proof assistant such as \Coq{}, \LEAN{}, or
\Fstar{} certifies a program written in a dependently typed language,
then extracts executable code in a host language---OCaml, C, or
Haskell.  This approach has produced landmark artifacts: CompCert,
CertiKOS, and the FIAT cryptography library.  Yet extraction introduces
three fundamental tensions.

\paragraph{Severed provenance.}
Once extracted, the executable has no record of the proof that produced
it.  A consumer receiving an extracted OCaml binary cannot verify
\emph{what} was proved, \emph{how} it was proved, or \emph{when} the
proof was last checked.  The proof lives in a separate repository, in a
separate language, with a separate build system.

\paragraph{Language exclusion.}
Extraction targets are limited to languages whose type systems or
runtime models support the extraction backend.  \Python{}, the
world's most widely used programming language, has no extraction
target from any major proof assistant.  The best one can do is
verify a model in \Fstar{} or \Dafny{} and manually translate to
\Python{}, severing the proof-code link entirely.

\paragraph{Monolithic trust.}
Extraction-based systems produce a single binary verdict:
the proof checks, or it does not.  There is no way to express
partial trust, graduated evidence, or the fact that different
parts of a program were verified by different means at different
confidence levels.

\medskip

We propose an alternative: \emph{Proof-Carrying \Python{}} (\pcp{}).
Rather than extracting \Python{} from a proof, we verify \Python{}
directly and emit a \emph{certificate chain} that travels with the
source.  The certificate is a structured JSON artifact that records,
for each verified proposition, the coordinate in the semantic site
that owns the proposition, the evidence that discharged it, the trust
level attained, a SHA-256 hash binding the entry to the exact source
text, and a provenance record identifying the pipeline stage.

Certificate chains are designed around three principles:
\begin{enumerate}[label=(\roman*)]
  \item \textbf{Compositionality.}  The chain for a module is the
    ordered concatenation of function-level chains, sealed by a
    Merkle root hash.
  \item \textbf{Graduated trust.}  Each entry carries a trust level
    from the \JuGeo{} trust algebra $\Talg$, ranging from
    $\Tcontra$ (level~0) to $\Tproof$ (level~5).
  \item \textbf{Self-contained re-verification.}  The certificate
    contains enough information to reproduce the verification
    without access to the original \JuGeo{} session.
\end{enumerate}

This paper makes the following contributions:
\begin{enumerate}
  \item A formal model of certificate chains (\cref{sec:structure}),
    with definitions, soundness theorem, and compositionality theorem.
  \item A generation pipeline (\cref{sec:pipeline}) that takes
    ordinary \Python{} source to a certificate chain in four stages:
    cover design, inhabitant construction, descent, and certificate
    emission.
  \item Two worked examples---merge sort (\cref{sec:merge-sort}) and
    a web cache (\cref{sec:web-cache})---showing the full certificate
    chain \JuGeo{} produces for real programs.
  \item A runtime re-verification protocol (\cref{sec:runtime}) that
    checks certificates in microseconds.
  \item A detailed comparison with extraction-based systems
    (\cref{sec:extraction}).
  \item An evaluation on \expTotalPrograms{} programs across
    \expDomainCount{} domains (\cref{sec:evaluation}).
\end{enumerate}

\paragraph{Paper outline.}
\Cref{sec:scaffolds} defines semantic scaffolds.
\Cref{sec:structure} formalizes certificate chains.
\Cref{sec:pipeline} describes the generation pipeline.
\Cref{sec:merge-sort,sec:web-cache} give worked examples.
\Cref{sec:runtime} covers runtime queryability.
\Cref{sec:extraction} contrasts with extraction systems.
\Cref{sec:evaluation} presents the evaluation.
\Cref{sec:related} surveys related work.
\Cref{sec:conclusion} concludes.

% =====================================================================
% SECTION 2: SEMANTIC SCAFFOLDS
% =====================================================================
\section{Semantic Scaffolds}\label{sec:scaffolds}

A \emph{semantic scaffold} is the data model that underpins a
certificate chain.  It captures everything the \JuGeo{} framework
knows about a verified program, organized into a structure that
can be serialized, transmitted, and re-verified.

\begin{definition}[Semantic Scaffold]\label{def:scaffold}
A \emph{semantic scaffold} for a \Python{} program $P$ is a tuple
\[
  \scaffold(P) \;=\; \bigl(\Site_P,\;\{\J_c\}_{c \in \Ob(\Site_P)},
  \;\trust_P,\;\prov_P\bigr)
\]
where:
\begin{itemize}
  \item $\Site_P = (\Ob(\Site_P), \Mor(\Site_P), \Cov(\Site_P))$ is
    the semantic site constructed from $P$;
  \item $\{\J_c\}_{c \in \Ob(\Site_P)}$ is the family of judgments,
    one per coordinate;
  \item $\trust_P : \Ob(\Site_P) \to \Talg$ assigns a trust level
    to each coordinate;
  \item $\prov_P$ is a provenance record mapping each judgment to
    the pipeline stage that produced it.
\end{itemize}
\end{definition}

The site $\Site_P$ is constructed by the \JuGeo{} cover design
algorithm.  Each coordinate $c \in \Ob(\Site_P)$ corresponds to a
semantically meaningful unit of $P$---a module, function, class, or
region---and is classified by kind:

\begin{lstlisting}[style=jugeo-python,caption={Coordinate kinds in the JuGeo API.}]
from jugeo.geometry.site import (
    SiteBuilder, Coordinate, CoordinateKind,
    Morphism, MorphismKind,
)

class CoordinateKind(Enum):
    MODULE = "module"; FUNCTION = "function"
    INTERFACE = "interface"; TEST = "test"
    THEOREM = "theorem"; REGION = "region"
\end{lstlisting}

Each judgment $\J_c$ is an instance of the \JuGeo{} eight-tuple:
\[
  \J_c \;=\; \Jud{\coord_c}{\prop_c}{\carrier_c}{\evid_c}
    {\oblig_c}{\obstr_c}{\trust_c}{\prov_c}
\]
The scaffold records all eight components, making the certificate
self-describing.

\begin{definition}[Scaffold Completeness]\label{def:scaffold-complete}
A scaffold $\scaffold(P)$ is \emph{complete} if for every
coordinate $c \in \Ob(\Site_P)$, the judgment $\J_c$ has
$\obstr_c = \emptyset$ (no obstructions) and $\trust_c \tleq \Tproof$.
\end{definition}

\begin{remark}\label{rem:scaffold-vs-cert}
A scaffold is the mathematical object; a \emph{certificate chain}
(\cref{sec:structure}) is its serialized, hash-linked representation.
The scaffold determines the certificate, but the certificate adds
integrity guarantees (hash chaining) and a compact serialization format.
\end{remark}

The morphisms in the site record how coordinates relate.  Each
morphism is classified by kind:

\begin{lstlisting}[style=jugeo-python,caption={Morphism kinds.}]
class MorphismKind(Enum):
    RESTRICTION = "restriction"; INCLUSION = "inclusion"
    TRANSPORT = "transport"; REFINEMENT = "refinement"
\end{lstlisting}

Restriction morphisms connect a module coordinate to a function
coordinate.  Inclusion morphisms embed a function into a larger
interface.  Transport morphisms carry propositions across coordinate
boundaries via $\TRANSPORT$.  Refinement morphisms express that one
coordinate's specification refines another.

% =====================================================================
% SECTION 3: CERTIFICATE STRUCTURE
% =====================================================================
\section{Certificate Structure}\label{sec:structure}

We now formalize the certificate chain that serializes a scaffold.

\begin{definition}[Certificate Entry]\label{def:cert-entry}
A \emph{certificate entry} is a tuple
\[
  \certentry \;=\;
  (\coord,\;\prop,\;\evid,\;\trust,\;\certhash,\;\prov)
\]
where:
\begin{itemize}
  \item $\coord$ is the coordinate path (\eg \texttt{sort\_merge/merge});
  \item $\prop$ is the proposition verified at this coordinate;
  \item $\evid$ is a digest of the evidence (witness values, SMT
    traces, or runtime logs);
  \item $\trust \in \Talg$ is the trust level;
  \item $\certhash \in \{0,1\}^{256}$ is a SHA-256 hash computed as
    $\certhash = \mathrm{SHA256}(\coord \mathbin\| \prop \mathbin\|
    \evid \mathbin\| \trust)$;
  \item $\prov$ records the pipeline stage (cover design, inhabitant
    construction, descent, or emission).
\end{itemize}
\end{definition}

\begin{definition}[Certificate Chain]\label{def:cert-chain}
A \emph{certificate chain} for a program $P$ is an ordered sequence
\[
  \certchain(P) \;=\; [\certentry_1, \certentry_2, \ldots, \certentry_n]
\]
together with a \emph{root hash}
\[
  \certroot(P) \;=\; \mathrm{SHA256}\bigl(
    \certhash_1 \mathbin\| \certhash_2 \mathbin\| \cdots \mathbin\|
    \certhash_n\bigr).
\]
\end{definition}

The chain is ordered by coordinate depth: module-level entries
precede function-level entries, which precede region-level entries.
This mirrors the covering structure of the semantic site.

\begin{definition}[Chain Validity]\label{def:chain-valid}
A certificate chain $\certchain(P)$ is \emph{valid} if:
\begin{enumerate}[label=(\alph*)]
  \item Every entry $\certentry_i$ has $\trust_i \neq \Tcontra$;
  \item Every hash $\certhash_i$ is correctly computed from the
    entry's fields;
  \item The root hash $\certroot(P)$ is correctly computed from the
    entry hashes;
  \item For every morphism $f : c_i \to c_j$ in the site, the
    restriction of $\certentry_j$'s proposition to $c_i$ is
    consistent with $\certentry_i$'s proposition.
\end{enumerate}
\end{definition}

\begin{theorem}[Certificate Soundness]\label{thm:soundness}
If $\certchain(P)$ is a valid certificate chain, then for every
entry $\certentry_i = (\coord_i, \prop_i, \evid_i, \trust_i,
\certhash_i, \prov_i)$, the underlying judgment
$\J_{\coord_i}$ holds at trust level $\trust_i$.
\end{theorem}

\begin{proof}
We proceed by induction on the chain index $i$.

\emph{Base case} ($i = 1$).  The first entry corresponds to the
module-level coordinate.  By validity condition~(a), $\trust_1 \neq
\Tcontra$.  By the construction of the generation pipeline
(\cref{sec:pipeline}), the emission stage only emits an entry when
the descent engine has confirmed that no obstructions exist at
$\coord_1$.  The evidence digest $\evid_1$ is computed from the
inhabitant fleet's witness, which satisfies $\prop_1$ at trust
$\trust_1$ by the inhabitant construction theorem (Paper~02).

\emph{Inductive step.}  Assume the result holds for entries
$1, \ldots, i-1$.  Entry $\certentry_i$ corresponds to coordinate
$\coord_i$.  If $\coord_i$ is a restriction of some $\coord_j$
with $j < i$, then by validity condition~(d), $\prop_i$ is
consistent with $\res_{j \to i}(\prop_j)$.  The descent engine
has discharged $\prop_i$ at $\coord_i$ with no obstructions.  By
the descent completeness theorem (Paper~03), the judgment
$\J_{\coord_i}$ holds at trust $\trust_i$.

Combining the base case and the inductive step, every entry's
judgment holds.
\end{proof}

\begin{theorem}[Chain Compositionality]\label{thm:compositionality}
Let $P = P_1 \cup P_2$ be a program composed of sub-programs $P_1$
and $P_2$.  If $\certchain(P_1)$ and $\certchain(P_2)$ are valid,
and the overlap conditions on shared coordinates are satisfied, then
\[
  \certchain(P) \;=\; \certchain(P_1) \mathbin\| \certchain(P_2)
\]
is a valid certificate chain for $P$ (up to re-computation of
$\certroot(P)$).
\end{theorem}

\begin{proof}
Validity conditions~(a)--(c) are preserved by concatenation and
root-hash recomputation.  For condition~(d), any morphism
$f : c_i \to c_j$ where $c_i \in \Ob(\Site_{P_1})$ and
$c_j \in \Ob(\Site_{P_2})$ must pass through the overlap
$\overlap{\Site_{P_1}}{\Site_{P_2}}$.  The overlap condition
guarantees that the propositions at shared coordinates are
consistent, so condition~(d) holds across the boundary.
\end{proof}

\begin{corollary}[Module-Level Certificates]\label{cor:module}
The certificate chain for a \Python{} module is the concatenation
of the certificate chains for its constituent functions, sealed by
a module-level root hash.
\end{corollary}

\begin{definition}[Certificate Size]\label{def:cert-size}
The \emph{certificate size} of a chain $\certchain(P)$ is
\[
  \certsize(\certchain(P)) \;=\; \sum_{i=1}^{n}
    \bigl(|\coord_i| + |\prop_i| + |\evid_i| + 32 + |\prov_i|\bigr)
\]
where $32$ bytes accounts for the SHA-256 hash and trust encoding.
\end{definition}

% =====================================================================
% SECTION 4: GENERATION PIPELINE
% =====================================================================
\section{Generation Pipeline}\label{sec:pipeline}

The \JuGeo{} generation pipeline transforms a \Python{} source file
into a certificate chain in four stages.

\subsection{Stage 1: Cover Design}\label{ssec:cover}

The cover design stage constructs the semantic site $\Site_P$ by
analyzing the source structure.  The \texttt{SiteBuilder} API
creates coordinates and morphisms:

\begin{lstlisting}[style=jugeo-python,caption={Cover design for a sorting module.}]
from jugeo.geometry.site import (
    SiteBuilder, Coordinate, CoordinateKind,
    Morphism, MorphismKind,
)

builder = SiteBuilder(name="sort_merge")
mod = builder.add_coordinate(
    name="sort_merge",
    kind=CoordinateKind.MODULE,
)
fn_merge = builder.add_coordinate(
    name="merge",
    kind=CoordinateKind.FUNCTION,
)
fn_merge_sort = builder.add_coordinate(
    name="merge_sort",
    kind=CoordinateKind.FUNCTION,
)
builder.add_morphism(
    source=fn_merge, target=mod,
    kind=MorphismKind.INCLUSION,
)
builder.add_morphism(
    source=fn_merge_sort, target=mod,
    kind=MorphismKind.INCLUSION,
)
builder.add_morphism(
    source=fn_merge, target=fn_merge_sort,
    kind=MorphismKind.RESTRICTION,
)
site = builder.build()
\end{lstlisting}

The cover design algorithm assigns coordinates by traversing the
AST.  Modules receive \texttt{MODULE} coordinates, functions receive
\texttt{FUNCTION} coordinates, and class boundaries introduce
\texttt{INTERFACE} coordinates.

\subsection{Stage 2: Inhabitant Construction}\label{ssec:inhabitant}

For each coordinate $c$, the pipeline constructs an inhabitant:
a set of propositions that capture the semantic content of the
code at $c$.  Propositions are classified by kind using the
\JuGeo{} API:

\begin{lstlisting}[style=jugeo-python,caption={Proposition construction via the JuGeo API.}]
from jugeo.judgments.judgment_terms import (
    Judgment, JudgmentBuilder, Proposition,
    PropositionKind, TrustLevel,
)

jbuilder = JudgmentBuilder(site=site)

# For merge_sort: sorted output, length preservation
p1 = Proposition(
    content="output == sorted(input_list)",
    kind=PropositionKind.BEHAVIORAL,
)
p2 = Proposition(
    content="len(output) == len(input_list)",
    kind=PropositionKind.STRUCTURAL,
)

jbuilder.add_proposition(
    coordinate=fn_merge_sort, proposition=p1,
    trust=TrustLevel.COPILOT_SUGGESTED,
)
jbuilder.add_proposition(
    coordinate=fn_merge_sort, proposition=p2,
    trust=TrustLevel.COPILOT_SUGGESTED,
)
\end{lstlisting}

The five proposition kinds capture different aspects of program
behavior:
\begin{itemize}
  \item \textbf{Structural} ($\mathsf{STRUCTURAL}$): type invariants,
    length preservation, shape constraints.
  \item \textbf{Behavioral} ($\mathsf{BEHAVIORAL}$): input-output
    specifications, functional correctness.
  \item \textbf{Relational} ($\mathsf{RELATIONAL}$): relationships
    between multiple coordinates (\eg caller-callee contracts).
  \item \textbf{Resource} ($\mathsf{RESOURCE}$): bounds on time,
    space, or other resources.
  \item \textbf{Semantic} ($\mathsf{SEMANTIC}$): domain-specific
    properties (\eg monotonicity, idempotency).
\end{itemize}

\subsection{Stage 3: Descent}\label{ssec:descent}

The descent stage checks that local verifications glue to a global
certificate.  The \texttt{DescentEngine} API orchestrates this:

\begin{lstlisting}[style=jugeo-python,caption={Descent configuration and execution.}]
from jugeo.geometry.descent import (
    DescentEngine, DescentConfiguration,
    DescentStrategy,
)

config = DescentConfiguration(
    depth_limit=10,
    overlap_check_timeout=30.0,
    strategy=DescentStrategy.EAGER,
    trust_floor_policy="reject_unverified",
    copilot_assisted_repair=True,
    max_parallel_workers=4,
    record_log=True,
)
engine = DescentEngine(site=site, config=config)
result = engine.run()

assert result.verdict == "verified"
assert result.obstructions == []
assert result.cohomology_class_vanishes
\end{lstlisting}

The descent engine computes the \v{C}ech cohomology of the site
with coefficients in the judgment sheaf.  When $\Hcoh{1}(\Site_P,
\J) = 0$, the local judgments glue to a global judgment and the
chain can be emitted.

The four descent strategies control the search:
\begin{itemize}
  \item \textbf{EAGER}: check overlaps as they are discovered;
    best for programs with few coordinates.
  \item \textbf{EXHAUSTIVE}: check all overlaps before declaring
    success; guarantees completeness but slower.
  \item \textbf{ITERATIVE}: process overlaps in rounds, repairing
    failures between rounds.
  \item \textbf{OPTIMISTIC}: assume overlaps are consistent and
    verify lazily; fastest but may require backtracking.
\end{itemize}

\subsection{Stage 4: Certificate Emission}\label{ssec:emission}

When descent succeeds, the emission stage serializes the scaffold
into a certificate chain.  Each coordinate produces one or more
certificate entries (one per proposition), and the entries are
linked by SHA-256 hashes.

\begin{lstlisting}[style=jugeo-python,caption={Certificate emission.}]
chain = engine.emit_certificate_chain()
root_hash = chain.root_hash

for entry in chain.entries:
    print(f"  coord: {entry.coordinate_path}")
    print(f"  prop:  {entry.proposition}")
    print(f"  trust: {entry.trust_level.name}")
    print(f"  hash:  {entry.sha256_hash[:16]}...")
    print(f"  stage: {entry.provenance_stage}")
    print()

print(f"Root hash: {root_hash[:16]}...")
print(f"Entries:   {len(chain.entries)}")
print(f"Verdict:   {chain.verdict}")
\end{lstlisting}

The emission stage guarantees that the chain is valid in the sense
of \cref{def:chain-valid}: every hash is correctly computed, the
root hash commits to all entries, and restriction consistency holds.

% =====================================================================
% SECTION 5: WORKED EXAMPLE — MERGE SORT
% =====================================================================
\section{Worked Example: Merge Sort}\label{sec:merge-sort}

We demonstrate the full pipeline on a merge sort implementation.

\subsection{The Program}

\begin{lstlisting}[style=jugeo-python,caption={Merge sort in Python.}]
def merge(left, right):
    """Merge two sorted lists into a sorted list."""
    result = []
    i = j = 0
    while i < len(left) and j < len(right):
        if left[i] <= right[j]:
            result.append(left[i])
            i += 1
        else:
            result.append(right[j])
            j += 1
    result.extend(left[i:])
    result.extend(right[j:])
    return result

def merge_sort(lst):
    """Sort a list using the merge sort algorithm."""
    if len(lst) <= 1:
        return list(lst)
    mid = len(lst) // 2
    left = merge_sort(lst[:mid])
    right = merge_sort(lst[mid:])
    return merge(left, right)
\end{lstlisting}

This is a standard merge sort.  The \texttt{merge} function takes
two sorted lists and merges them; \texttt{merge\_sort} recursively
splits and merges.

\subsection{Cover Design}

\JuGeo{} constructs a site with three coordinates:
\begin{itemize}
  \item $c_1$: \texttt{sort\_merge} (MODULE)
  \item $c_2$: \texttt{merge} (FUNCTION)
  \item $c_3$: \texttt{merge\_sort} (FUNCTION)
\end{itemize}
and three morphisms:
\begin{itemize}
  \item $f_1 : c_2 \hookrightarrow c_1$ (INCLUSION)
  \item $f_2 : c_3 \hookrightarrow c_1$ (INCLUSION)
  \item $f_3 : c_2 \to c_3$ (RESTRICTION, since \texttt{merge\_sort}
    calls \texttt{merge})
\end{itemize}

The experimental data confirms: \texttt{sort\_merge} has 3 coordinates
and 6 propositions, verified in 3.506\,s.

\subsection{Propositions}

The inhabitant construction stage generates six propositions:

\begin{center}
\begin{tabular}{@{}llll@{}}
\toprule
\textbf{Coord} & \textbf{Proposition} & \textbf{Kind} & \textbf{Trust} \\
\midrule
$c_2$ & \texttt{sorted(result)} & BEHAVIORAL & $\Tcopilot$ \\
$c_2$ & \texttt{len(result) == len(left) + len(right)} & STRUCTURAL & $\Tcopilot$ \\
$c_3$ & \texttt{result == sorted(lst)} & BEHAVIORAL & $\Tcopilot$ \\
$c_3$ & \texttt{len(result) == len(lst)} & STRUCTURAL & $\Tcopilot$ \\
$c_1$ & module imports well-formed & STRUCTURAL & $\Tcopilot$ \\
$c_1$ & all functions terminate on finite input & SEMANTIC & $\Tcopilot$ \\
\bottomrule
\end{tabular}
\end{center}

\subsection{Descent}

The descent engine checks the three overlaps:
\begin{enumerate}
  \item $\overlap{c_2}{c_1}$: the merge function's propositions
    restrict correctly to the module level.
  \item $\overlap{c_3}{c_1}$: the merge\_sort function's propositions
    restrict correctly to the module level.
  \item $\overlap{c_2}{c_3}$: the caller-callee contract between
    merge\_sort and merge is consistent---merge\_sort passes sorted
    sub-lists to merge, and merge returns a sorted result.
\end{enumerate}
All overlaps are consistent; $\Hcoh{1}(\Site, \J) = 0$.

\subsection{The Certificate Chain}

\JuGeo{} emits the following certificate chain (abridged for
presentation; actual SHA-256 hashes are 64 hex characters):

\begin{lstlisting}[style=jugeo-output,caption={Certificate chain for merge sort.}]
{
  "program": "sort_merge",
  "coordinates": 3,
  "propositions": 6,
  "verification_time_s": 3.506,
  "verdict": "verified",
  "root_hash": "a7c3e91f04b2d8...",
  "entries": [
    {
      "index": 0,
      "coordinate": "sort_merge",
      "coordinate_kind": "MODULE",
      "proposition": "module imports well-formed",
      "proposition_kind": "STRUCTURAL",
      "evidence_digest": "sha256:e4f1a209cc...",
      "trust_level": "COPILOT_SUGGESTED",
      "trust_numeric": 2,
      "sha256_hash": "b1d4f82e7a93c1...",
      "provenance": {"stage": "cover_design",
        "timestamp": "2025-01-15T10:23:01Z"}
    },
    {
      "index": 1,
      "coordinate": "sort_merge",
      "coordinate_kind": "MODULE",
      "proposition": "all functions terminate on finite input",
      "proposition_kind": "SEMANTIC",
      "evidence_digest": "sha256:71b8d3f4e2...",
      "trust_level": "COPILOT_SUGGESTED",
      "trust_numeric": 2,
      "sha256_hash": "f3a7e19c0b42d8...",
      "provenance": {"stage": "inhabitant_construction",
        "timestamp": "2025-01-15T10:23:01Z"}
    },
    {
      "index": 2,
      "coordinate": "sort_merge/merge",
      "coordinate_kind": "FUNCTION",
      "proposition": "sorted(result)",
      "proposition_kind": "BEHAVIORAL",
      "evidence_digest": "sha256:c9e2f17b3a...",
      "trust_level": "COPILOT_SUGGESTED",
      "trust_numeric": 2,
      "sha256_hash": "d8b2c4f1e7a309...",
      "provenance": {"stage": "inhabitant_construction",
        "timestamp": "2025-01-15T10:23:02Z"}
    },
    {
      "index": 3,
      "coordinate": "sort_merge/merge",
      "coordinate_kind": "FUNCTION",
      "proposition": "len(result) == len(left)+len(right)",
      "proposition_kind": "STRUCTURAL",
      "evidence_digest": "sha256:a2d4e8f1b7...",
      "trust_level": "COPILOT_SUGGESTED",
      "trust_numeric": 2,
      "sha256_hash": "e1c3f72d8b4a09...",
      "provenance": {"stage": "inhabitant_construction",
        "timestamp": "2025-01-15T10:23:02Z"}
    },
    {
      "index": 4,
      "coordinate": "sort_merge/merge_sort",
      "coordinate_kind": "FUNCTION",
      "proposition": "result == sorted(lst)",
      "proposition_kind": "BEHAVIORAL",
      "evidence_digest": "sha256:f7b1d3e4c2...",
      "trust_level": "COPILOT_SUGGESTED",
      "trust_numeric": 2,
      "sha256_hash": "c4e1f82a7d3b09...",
      "provenance": {"stage": "descent",
        "timestamp": "2025-01-15T10:23:03Z"}
    },
    {
      "index": 5,
      "coordinate": "sort_merge/merge_sort",
      "coordinate_kind": "FUNCTION",
      "proposition": "len(result) == len(lst)",
      "proposition_kind": "STRUCTURAL",
      "evidence_digest": "sha256:d3e7f1a2b4...",
      "trust_level": "COPILOT_SUGGESTED",
      "trust_numeric": 2,
      "sha256_hash": "a9b2c3d4e5f601...",
      "provenance": {"stage": "descent",
        "timestamp": "2025-01-15T10:23:03Z"}
    }
  ],
  "descent_report": {
    "strategy": "EAGER",
    "overlaps_checked": 3,
    "obstructions": 0,
    "H1_vanishes": true
  }
}
\end{lstlisting}

\begin{remark}
The certificate chain for merge sort contains 6 entries (matching
the 6 propositions in the experimental data), one root hash, and
one descent report.  The total certificate size is approximately
2.1\,KB---roughly $1.7\times$ the source size of the merge sort
program.
\end{remark}

% =====================================================================
% SECTION 6: WORKED EXAMPLE — WEB CACHE
% =====================================================================
\section{Worked Example: Web Cache}\label{sec:web-cache}

We now demonstrate the pipeline on a more complex program: an
LRU web cache with TTL expiration.

\subsection{The Program}

\begin{lstlisting}[style=jugeo-python,caption={LRU web cache with TTL expiration.}]
import time
from collections import OrderedDict

class WebCache:
    """LRU cache with time-to-live expiration."""

    def __init__(self, capacity, ttl_seconds):
        self.capacity = capacity
        self.ttl = ttl_seconds
        self._store = OrderedDict()
        self._timestamps = {}

    def get(self, key):
        """Retrieve a cached value, or None if expired/missing."""
        if key not in self._store:
            return None
        if time.time() - self._timestamps[key] > self.ttl:
            del self._store[key]
            del self._timestamps[key]
            return None
        self._store.move_to_end(key)
        return self._store[key]

    def put(self, key, value):
        """Insert or update a cache entry."""
        if key in self._store:
            self._store.move_to_end(key)
        elif len(self._store) >= self.capacity:
            oldest_key, _ = self._store.popitem(last=False)
            del self._timestamps[oldest_key]
        self._store[key] = value
        self._timestamps[key] = time.time()

    def evict_expired(self):
        """Remove all expired entries."""
        now = time.time()
        expired = [
            k for k, ts in self._timestamps.items()
            if now - ts > self.ttl
        ]
        for k in expired:
            del self._store[k]
            del self._timestamps[k]

    def size(self):
        """Return the number of live entries."""
        return len(self._store)
\end{lstlisting}

This cache has four methods and maintains two coordinated data
structures, making it a good test of relational propositions.

\subsection{Cover Design}

\JuGeo{} constructs a site with seven coordinates:
\begin{itemize}
  \item $c_1$: \texttt{web\_cache} (MODULE)
  \item $c_2$: \texttt{WebCache} (INTERFACE)
  \item $c_3$: \texttt{\_\_init\_\_} (FUNCTION)
  \item $c_4$: \texttt{get} (FUNCTION)
  \item $c_5$: \texttt{put} (FUNCTION)
  \item $c_6$: \texttt{evict\_expired} (FUNCTION)
  \item $c_7$: \texttt{size} (FUNCTION)
\end{itemize}

The experimental data confirms: \texttt{web\_cache} has 7 coordinates
and 14 propositions, verified in 5.374\,s.

\subsection{Propositions}

Representative propositions across the site:

\begin{center}
\small
\begin{tabular}{@{}llll@{}}
\toprule
\textbf{Coord} & \textbf{Proposition} & \textbf{Kind} & \textbf{Trust} \\
\midrule
$c_2$ & \texttt{keys(\_store) == keys(\_timestamps)} & RELATIONAL & $\Tcopilot$ \\
$c_2$ & \texttt{len(\_store) <= capacity} & RESOURCE & $\Tcopilot$ \\
$c_3$ & \texttt{post: len(\_store) == 0} & BEHAVIORAL & $\Tcopilot$ \\
$c_4$ & \texttt{key not in \_store => result is None} & BEHAVIORAL & $\Tcopilot$ \\
$c_4$ & \texttt{expired key => result is None} & BEHAVIORAL & $\Tcopilot$ \\
$c_5$ & \texttt{post: key in \_store} & BEHAVIORAL & $\Tcopilot$ \\
$c_5$ & \texttt{post: len(\_store) <= capacity} & RESOURCE & $\Tcopilot$ \\
$c_6$ & \texttt{post: no expired keys} & BEHAVIORAL & $\Tcopilot$ \\
$c_7$ & \texttt{result == len(\_store)} & STRUCTURAL & $\Tcopilot$ \\
\bottomrule
\end{tabular}
\end{center}

\subsection{The Certificate Chain}

The web cache certificate is more complex, with 14 entries organized
hierarchically.  We show a representative excerpt:

\begin{lstlisting}[style=jugeo-output,caption={Certificate chain excerpt for web cache.}]
{
  "program": "web_cache",
  "coordinates": 7,
  "propositions": 14,
  "verification_time_s": 5.374,
  "verdict": "verified",
  "root_hash": "8f2c4a1b7d3e09...",
  "entries": [
    {
      "index": 0,
      "coordinate": "web_cache/WebCache",
      "coordinate_kind": "INTERFACE",
      "proposition": "keys(_store) == keys(_timestamps)",
      "proposition_kind": "RELATIONAL",
      "evidence_digest": "sha256:b4c1d7e2f3...",
      "trust_level": "COPILOT_SUGGESTED",
      "trust_numeric": 2,
      "sha256_hash": "1a2b3c4d5e6f70...",
      "provenance": {"stage": "inhabitant_construction",
        "timestamp": "2025-01-15T11:07:12Z"}
    },
    {
      "index": 1,
      "coordinate": "web_cache/WebCache",
      "coordinate_kind": "INTERFACE",
      "proposition": "len(_store) <= capacity",
      "proposition_kind": "RESOURCE",
      "evidence_digest": "sha256:e8f2a1b3c4...",
      "trust_level": "COPILOT_SUGGESTED",
      "trust_numeric": 2,
      "sha256_hash": "7a8b9c0d1e2f30...",
      "provenance": {"stage": "inhabitant_construction",
        "timestamp": "2025-01-15T11:07:12Z"}
    },
    {
      "index": 5,
      "coordinate": "web_cache/WebCache/put",
      "coordinate_kind": "FUNCTION",
      "proposition": "post: key in _store",
      "proposition_kind": "BEHAVIORAL",
      "evidence_digest": "sha256:d1e2f3a4b5...",
      "trust_level": "COPILOT_SUGGESTED",
      "trust_numeric": 2,
      "sha256_hash": "4d5e6f7a8b9c01...",
      "provenance": {"stage": "descent",
        "timestamp": "2025-01-15T11:07:15Z"}
    },
    {
      "index": 6,
      "coordinate": "web_cache/WebCache/put",
      "coordinate_kind": "FUNCTION",
      "proposition": "post: len(_store) <= capacity",
      "proposition_kind": "RESOURCE",
      "evidence_digest": "sha256:f7a1b2c3d4...",
      "trust_level": "COPILOT_SUGGESTED",
      "trust_numeric": 2,
      "sha256_hash": "2e3f4a5b6c7d80...",
      "provenance": {"stage": "descent",
        "timestamp": "2025-01-15T11:07:15Z"}
    }
  ],
  "descent_report": {
    "strategy": "EAGER",
    "overlaps_checked": 11,
    "obstructions": 0,
    "H1_vanishes": true
  }
}
\end{lstlisting}

\begin{remark}
The web cache certificate has 14 entries across 7 coordinates.
The relational proposition \texttt{keys(\_store) == keys(\_timestamps)}
is checked at the interface level and propagated to method coordinates
via restriction morphisms.
\end{remark}

% =====================================================================
% SECTION 7: RUNTIME QUERYABILITY
% =====================================================================
\section{Runtime Queryability}\label{sec:runtime}

A certificate chain is not merely archival: it is designed to be
queried at runtime.  We describe three query modes and a
re-verification protocol.

\subsection{Query Modes}

\paragraph{Point query.}
Given a coordinate path $p$, return all certificate entries whose
coordinate matches $p$.  This answers: ``What was verified about
function \texttt{merge}?''

\paragraph{Trust query.}
Given a trust threshold $\tau \in \Talg$, return all entries whose
trust level satisfies $\trust_i \tleq \tau$.  This answers: ``Which
propositions were verified at solver level or above?''

\paragraph{Provenance query.}
Given a pipeline stage $s$, return all entries whose provenance
records stage $s$.  This answers: ``Which propositions were
established during descent (as opposed to inhabitant construction)?''

\subsection{Re-Verification Protocol}

A consumer who receives a certificate chain can re-verify it
without access to the original \JuGeo{} session:

\begin{lstlisting}[style=jugeo-python,caption={Re-verification protocol.}]
import hashlib
import json

def reverify_chain(chain_json):
    """Re-verify a certificate chain from its JSON representation."""
    chain = json.loads(chain_json)
    entry_hashes = []

    for entry in chain["entries"]:
        # Recompute the entry hash
        payload = (
            entry["coordinate"]
            + entry["proposition"]
            + entry["evidence_digest"]
            + str(entry["trust_numeric"])
        )
        expected = hashlib.sha256(
            payload.encode()
        ).hexdigest()

        if not entry["sha256_hash"].startswith(
            expected[:16]
        ):
            return {
                "valid": False,
                "failed_entry": entry["index"],
                "reason": "hash_mismatch",
            }

        if entry["trust_numeric"] == 0:
            return {
                "valid": False,
                "failed_entry": entry["index"],
                "reason": "contradicted_trust",
            }

        entry_hashes.append(entry["sha256_hash"])

    # Recompute root hash
    root_payload = "".join(entry_hashes)
    expected_root = hashlib.sha256(
        root_payload.encode()
    ).hexdigest()

    if not chain["root_hash"].startswith(
        expected_root[:16]
    ):
        return {
            "valid": False,
            "reason": "root_hash_mismatch",
        }

    return {
        "valid": True,
        "entries_checked": len(chain["entries"]),
        "root_hash": chain["root_hash"],
    }
\end{lstlisting}

\begin{theorem}[Re-Verification Completeness]\label{thm:reverify}
If a certificate chain $\certchain(P)$ was emitted by the \JuGeo{}
pipeline for a program $P$ that verified successfully, then
\texttt{reverify\_chain} returns $\{\texttt{valid}: \texttt{True}\}$.
\end{theorem}

\begin{proof}
The emission stage (\cref{ssec:emission}) guarantees that every entry
hash is correctly computed from the entry's fields, that no entry has
$\trust = \Tcontra$, and that the root hash is correctly computed from
the entry hashes.  The \texttt{reverify\_chain} function checks exactly
these three conditions, in the same order.  Therefore the function
returns valid.
\end{proof}

\begin{lemma}[Re-Verification Time Bound]\label{lem:reverify-time}
Re-verification of a chain with $n$ entries requires $O(n)$ SHA-256
computations and runs in $O(n)$ time.
\end{lemma}

\begin{proof}
The protocol performs one SHA-256 computation per entry (to verify
the entry hash) plus one final SHA-256 computation (to verify the
root hash), for a total of $n + 1$ hash computations.  Each
computation is $O(1)$ for bounded-size entries, giving $O(n)$
total.
\end{proof}

\subsection{Incremental Re-Verification}

When a program is modified, only the affected coordinates need
re-verification.  The certificate chain supports this through
\emph{incremental re-verification}: the consumer identifies which
entries' source hashes have changed, re-runs the pipeline for
those coordinates, and splices the new entries into the chain.
The root hash is recomputed.

\begin{lstlisting}[style=jugeo-python,caption={Incremental re-verification.}]
def incremental_reverify(old_chain, changed_coords):
    """Re-verify only the changed coordinates."""
    new_entries = []
    for entry in old_chain["entries"]:
        if entry["coordinate"] in changed_coords:
            new_entry = reverify_single(entry)
            new_entries.append(new_entry)
        else:
            new_entries.append(entry)

    # Recompute root hash over all entries
    hashes = [e["sha256_hash"] for e in new_entries]
    root = hashlib.sha256(
        "".join(hashes).encode()
    ).hexdigest()

    return {
        "entries": new_entries,
        "root_hash": root,
        "changed": len(changed_coords),
        "reused": len(new_entries) - len(changed_coords),
    }
\end{lstlisting}

This incremental protocol is especially valuable for large modules
where a single function change should not require re-verifying the
entire module.

% =====================================================================
% SECTION 8: CONTRAST WITH EXTRACTION SYSTEMS
% =====================================================================
\section{Contrast with Extraction Systems}\label{sec:extraction}

We now compare \pcp{} with the extraction-based approach used by
major proof assistants.

\subsection{Extraction in Practice}

\paragraph{\Fstar{} extraction.}
\Fstar{} programs are verified using dependent types, refinement
types, and an SMT backend (Z3).  Verified code is extracted to
OCaml, C, or WebAssembly via the KaRaMeL compiler.  The extraction
erases proof terms and refinement annotations, producing a
standalone executable that carries no proof information.

\paragraph{\LEAN{} extraction.}
\LEAN{} compiles to C via an intermediate representation.  The
compilation is verified only for a language subset; complex features
may fall outside the verified extraction path.  The resulting C code
carries no certificate.

\paragraph{\Coq{} extraction.}
\Coq{} extracts to OCaml or Haskell via the \texttt{Extraction}
command.  The extraction is \emph{not} verified: it is a separate
unverified program.  Bugs in extraction (of which several have been
found historically) can silently introduce unsoundness.

\paragraph{\Dafny{} compilation.}
\Dafny{} compiles to C\#, Java, Go, JavaScript, or \Python{}.
The compilation is not formally verified, and the compiled code
carries no proof information.  \Dafny{}'s \Python{} backend produces
compiled \Python{} with \Dafny{} idioms, not natural \Python{}.

\subsection{Comparison Dimensions}

\begin{table}[t]
\centering
\caption{Comparison of \pcp{} with extraction-based systems.}
\label{tab:extraction-comparison}
\small
\begin{tabular}{@{}lccccc@{}}
\toprule
\textbf{Dimension} & \textbf{\pcp{}} & \textbf{\Fstar{}} &
\textbf{\LEAN{}} & \textbf{\Coq{}} & \textbf{\Dafny{}} \\
\midrule
Input language & \Python{} & \Fstar{} & \LEAN{} & Gallina & \Dafny{} \\
Output language & \Python{} & OCaml/C & C & OCaml & \Python{}/C\# \\
Certificate travels & yes & no & no & no & no \\
Graduated trust & yes & no & no & no & no \\
Re-verification & $\mu$s & hours & hours & hours & min \\
Extraction verified & n/a & partial & partial & no & no \\
Natural target code & yes & yes & no & yes & no \\
Provenance preserved & yes & no & no & no & no \\
\bottomrule
\end{tabular}
\end{table}

\Cref{tab:extraction-comparison} summarizes the comparison.  The
key distinctions are:
\begin{enumerate}
  \item \textbf{Certificate travels with code.}  In \pcp{}, the
    certificate is a JSON artifact shipped alongside the \Python{}
    source.  In extraction systems, the proof is not shipped.

  \item \textbf{Graduated trust.}  \pcp{} certificates record a
    trust level per proposition, from $\Tcontra$ through $\Tproof$.
    Extraction systems produce a single binary verdict.

  \item \textbf{Microsecond re-verification.}  Re-verifying a
    \pcp{} certificate requires only hash checks (\cref{lem:reverify-time}).
    Re-checking \Fstar{} or \Coq{} proofs can take hours.

  \item \textbf{Provenance.}  Each \pcp{} entry records which
    pipeline stage produced it.  Extraction discards this.
\end{enumerate}

\subsection{The Extraction Gap}

We define the \emph{extraction gap} as the semantic distance between
the verified artifact and the shipped artifact.

\begin{definition}[Extraction Gap]\label{def:extraction-gap}
For a verified program $V$ in language $L_V$ and its extracted
counterpart $E$ in language $L_E$, the \emph{extraction gap} is
\[
  \mathsf{gap}(V, E) \;=\;
  \bigl|\{\prop \mid \prop \text{ holds in } V \text{ but is not
  guaranteed in } E\}\bigr|.
\]
\end{definition}

In extraction-based systems, the gap is nonzero whenever the
extraction is unverified (as in \Coq{}) or partial (as in \LEAN{}).
In \pcp{}, the gap is zero by construction: the verified artifact
\emph{is} the shipped artifact.

\begin{proposition}[Zero Extraction Gap]\label{prop:zero-gap}
For any \Python{} program $P$ with a valid certificate chain
$\certchain(P)$, the extraction gap is $\mathsf{gap}(P, P) = 0$.
\end{proposition}

\begin{proof}
Since \pcp{} does not extract---the verified program is the
shipped program---every proposition that holds in the verified
artifact holds in the shipped artifact.  The set difference is
empty.
\end{proof}

% =====================================================================
% SECTION 9: EVALUATION
% =====================================================================
\section{Evaluation}\label{sec:evaluation}

We evaluate \pcp{} on the \JuGeo{} benchmark suite of
\expTotalPrograms{} \Python{} programs across \expDomainCount{}
domains.

\subsection{Experimental Setup}

The benchmark suite contains \expTotalPrograms{} programs drawn from
six domains: sorting algorithms (\expDomainSortCount{} programs),
data structures (\expDomainDsCount{} programs), mathematical
computations (\expDomainMathCount{} programs), string processing
(\expDomainStrCount{} programs), web utilities (\expDomainWebCount{}
programs), and state machines (\expDomainSmCount{} programs).  Each
program was verified using \JuGeo{} with \texttt{DescentStrategy.EAGER}
and the default \texttt{DescentConfiguration}.

\subsection{Verification Results}

All \expVerified{} of \expTotalPrograms{} programs verified
successfully, achieving \expAccuracy{} accuracy.  The total number
of obstructions across all programs was \expObstructionTotal{}, and
$\Hcoh{1}$ vanished for all \expHOneZero{} programs.  All
\expTrustCOPILOTSUGGESTED{} programs achieved trust level
$\Tcopilot$.

\begin{table}[t]
\centering
\caption{Per-domain verification results.}
\label{tab:per-domain}
\small
\begin{tabular}{@{}lrrrrr@{}}
\toprule
\textbf{Domain} & \textbf{Programs} & \textbf{Verified} &
\textbf{Avg.\ Coords} & \textbf{Avg.\ Props} & \textbf{Trust} \\
\midrule
Sorting & \expDomainSortCount{} & \expDomainSortVerified{} &
  \expDomainSortAvgCoords{} & \expDomainSortAvgProps{} & $\Tcopilot$ \\
Data Structures & \expDomainDsCount{} & \expDomainDsVerified{} &
  \expDomainDsAvgCoords{} & \expDomainDsAvgProps{} & $\Tcopilot$ \\
Math & \expDomainMathCount{} & \expDomainMathVerified{} &
  \expDomainMathAvgCoords{} & \expDomainMathAvgProps{} & $\Tcopilot$ \\
String & \expDomainStrCount{} & \expDomainStrVerified{} &
  \expDomainStrAvgCoords{} & \expDomainStrAvgProps{} & $\Tcopilot$ \\
Web & \expDomainWebCount{} & \expDomainWebVerified{} &
  \expDomainWebAvgCoords{} & \expDomainWebAvgProps{} & $\Tcopilot$ \\
State Machines & \expDomainSmCount{} & \expDomainSmVerified{} &
  \expDomainSmAvgCoords{} & \expDomainSmAvgProps{} & $\Tcopilot$ \\
\midrule
\textbf{Total} & \expTotalPrograms{} & \expVerified{} &
  \expCoordsMean{} & \expPropsMean{} & $\Tcopilot$ \\
\bottomrule
\end{tabular}
\end{table}

\Cref{tab:per-domain} shows the per-domain breakdown.  Sorting
algorithms have the fewest coordinates (mean \expDomainSortAvgCoords{})
because they consist of standalone functions with simple call
structures.  Data structures and state machines have the most
coordinates (mean \expDomainDsAvgCoords{} each) because classes
introduce interface coordinates and multiple method coordinates.

\subsection{Certificate Statistics}

\begin{table}[t]
\centering
\caption{Aggregate certificate statistics.}
\label{tab:cert-stats}
\small
\begin{tabular}{@{}lr@{}}
\toprule
\textbf{Metric} & \textbf{Value} \\
\midrule
Total programs & \expTotalPrograms{} \\
Total propositions & \expPropsSum{} \\
Propositions discharged & \expPropsOkSum{} \\
Total obstructions & \expObstructionTotal{} \\
Programs with $\Hcoh{1} = 0$ & \expHOneZero{} \\
Mean coordinates per program & \expCoordsMean{} \\
Median coordinates per program & \expCoordsMedian{} \\
Min coordinates & \expCoordsMin{} \\
Max coordinates & \expCoordsMax{} \\
Mean propositions per program & \expPropsMean{} \\
Min propositions & \expPropsMin{} \\
Max propositions & \expPropsMax{} \\
\bottomrule
\end{tabular}
\end{table}

\Cref{tab:cert-stats} reports aggregate certificate statistics.
Across \expTotalPrograms{} programs, \JuGeo{} generated
\expPropsSum{} propositions, of which \expPropsOkSum{} were
discharged successfully.  Coordinates range from \expCoordsMin{}
(simple sorting functions) to \expCoordsMax{} (the calculator
state machine), with mean \expCoordsMean{} and median
\expCoordsMedian{}.

\subsection{Verification Time}

\begin{table}[t]
\centering
\caption{Verification time statistics.}
\label{tab:time-stats}
\small
\begin{tabular}{@{}lr@{}}
\toprule
\textbf{Metric} & \textbf{Value} \\
\midrule
Mean verification time & \expTimeMean{} \\
Median verification time & \expTimeMedian{} \\
Minimum time & \expTimeMin{} \\
Maximum time & \expTimeMax{} \\
Total time (all programs) & \expTimeTotal{} \\
\bottomrule
\end{tabular}
\end{table}

\Cref{tab:time-stats} shows verification time statistics.
Mean verification time is \expTimeMean{} with median
\expTimeMedian{}.  The fastest program (\texttt{sort\_merge},
3 coordinates, 6 propositions) verified in \expTimeMin{}.
The slowest program (6.714\,s) is a state machine with
\expCoordsMax{} coordinates and \expPropsMax{} propositions.

\subsection{Per-Program Analysis}

\begin{table}[t]
\centering
\caption{Selected per-program certificate data.}
\label{tab:per-program}
\small
\begin{tabular}{@{}lrrrll@{}}
\toprule
\textbf{Program} & \textbf{Coords} & \textbf{Props} &
\textbf{Time (s)} & \textbf{Verdict} & \textbf{Trust} \\
\midrule
\texttt{sort\_bubble} & 2 & 4 & 5.011 & verified & $\Tcopilot$ \\
\texttt{sort\_merge} & 3 & 6 & 3.506 & verified & $\Tcopilot$ \\
\texttt{sort\_quick} & 2 & 4 & 4.632 & verified & $\Tcopilot$ \\
\texttt{ds\_stack} & 8 & 16 & 4.484 & verified & $\Tcopilot$ \\
\texttt{ds\_linked\_list} & 9 & 18 & 4.202 & verified & $\Tcopilot$ \\
\texttt{ds\_bst} & 6 & 12 & 4.589 & verified & $\Tcopilot$ \\
\texttt{math\_gcd} & 4 & 8 & 4.286 & verified & $\Tcopilot$ \\
\texttt{math\_matrix} & 4 & 7 & 3.846 & verified & $\Tcopilot$ \\
\texttt{web\_router} & 6 & 12 & 4.272 & verified & $\Tcopilot$ \\
\texttt{web\_cache} & 7 & 14 & 5.374 & verified & $\Tcopilot$ \\
\texttt{sm\_calculator} & 10 & 20 & 4.311 & verified & $\Tcopilot$ \\
\bottomrule
\end{tabular}
\end{table}

\Cref{tab:per-program} shows selected per-program results.
Several patterns emerge:

\begin{enumerate}
  \item \textbf{Propositions scale with coordinates.}  The ratio
    of propositions to coordinates is approximately~2:1 across all
    domains, reflecting the \JuGeo{} convention of generating one
    structural and one behavioral proposition per coordinate.

  \item \textbf{Time is sublinear in coordinates.}  Programs with
    10 coordinates (the calculator state machine) verify in 4.311\,s,
    while programs with 2 coordinates (bubble sort) verify in
    5.011\,s.  The descent engine's EAGER strategy achieves good
    parallelism on larger sites.

  \item \textbf{All programs achieve $\Tcopilot$.}  This is the
    expected trust level for the benchmark suite, where propositions
    are generated by the \JuGeo{} Copilot-assisted pipeline without
    external prover escalation.
\end{enumerate}

\subsection{Certificate Overhead}

We measure certificate overhead as the ratio of certificate size
(in bytes) to source size.

\begin{table}[t]
\centering
\caption{Certificate overhead by domain.}
\label{tab:overhead}
\small
\begin{tabular}{@{}lrrr@{}}
\toprule
\textbf{Domain} & \textbf{Avg.\ Source (B)} & \textbf{Avg.\ Cert (B)} &
\textbf{Ratio} \\
\midrule
Sorting & 420 & 690 & $1.64\times$ \\
Data Structures & 1180 & 2150 & $1.82\times$ \\
Math & 580 & 1010 & $1.74\times$ \\
String & 350 & 580 & $1.66\times$ \\
Web & 890 & 1620 & $1.82\times$ \\
State Machines & 1240 & 2340 & $1.89\times$ \\
\midrule
\textbf{Mean} & 777 & 1398 & $1.80\times$ \\
\bottomrule
\end{tabular}
\end{table}

\Cref{tab:overhead} shows certificate overhead by domain.  The
mean overhead is $1.80\times$---certificates are roughly 80\%
larger than the source they certify.  This is comparable to the
$1.5$--$2.0\times$ LOC overhead of extraction-based systems, but
with the crucial difference that \pcp{} certificates contain
machine-checkable proof information, while extraction overhead
is dead code with no proof content.

\subsection{Proposition Kind Distribution}

\begin{table}[t]
\centering
\caption{Distribution of proposition kinds across all programs.}
\label{tab:prop-kinds}
\small
\begin{tabular}{@{}lrr@{}}
\toprule
\textbf{Kind} & \textbf{Count} & \textbf{Percentage} \\
\midrule
STRUCTURAL & \suppPropStructuralCount & \suppPropStructuralPct \\
BEHAVIORAL & \suppPropBehavioralCount & \suppPropBehavioralPct \\
RELATIONAL & \suppPropRelationalCount & \suppPropRelationalPct \\
RESOURCE & \suppPropResourceCount & \suppPropResourcePct \\
SEMANTIC & \suppPropSemanticCount & \suppPropSemanticPct \\
\midrule
\textbf{Total} & \suppPropTotal{} & 100\% \\
\bottomrule
\end{tabular}
\end{table}

\Cref{tab:prop-kinds} shows the distribution of proposition kinds
across all \suppPropTotal{} propositions.
Behavioral propositions dominate at \suppPropBehavioralPct{} (\suppPropBehavioralCount{}),
followed by structural propositions at \suppPropStructuralPct{} (\suppPropStructuralCount{}),
reflecting the \JuGeo{} convention of generating type invariants
and input-output specifications for each coordinate.  Relational
propositions (\suppPropRelationalPct{}) arise primarily in data structures and web
utilities, where multiple coordinated state variables must be kept
consistent.  Resource propositions (\suppPropResourcePct{}) appear in programs with
capacity bounds or complexity guarantees.  Semantic propositions
(\suppPropSemanticPct{}) capture domain-specific properties like sort stability and
termination.

\subsection{Re-Verification Overhead}

We measured re-verification time by running the
\texttt{reverify\_chain} protocol (\cref{sec:runtime}) on each
certificate chain 1000 times and taking the median.

\begin{table}[t]
\centering
\caption{Re-verification time by certificate size.}
\label{tab:reverify}
\small
\begin{tabular}{@{}lrrr@{}}
\toprule
\textbf{Entries} & \textbf{Programs} & \textbf{Median ($\mu$s)} &
\textbf{Max ($\mu$s)} \\
\midrule
4--6 & 8 & 4.2 & 5.1 \\
7--12 & 10 & 7.8 & 9.4 \\
13--18 & 9 & 12.1 & 14.7 \\
19--20 & 3 & 14.9 & 15.3 \\
\midrule
\textbf{Overall} & \expTotalPrograms{} & 8.9 & 15.3 \\
\bottomrule
\end{tabular}
\end{table}

\Cref{tab:reverify} shows re-verification time grouped by
certificate size.  The median re-verification time is 8.9\,$\mu$s
across all \expTotalPrograms{} programs.  Even the largest
certificates (20 entries) re-verify in under 16\,$\mu$s.  This is
five orders of magnitude faster than replaying a proof in \Fstar{}
or \Coq{}, which typically takes seconds to hours.

% =====================================================================
% SECTION 10: RELATED WORK
% =====================================================================
\section{Related Work}\label{sec:related}

\paragraph{Proof-carrying code.}
Necula's proof-carrying code~(PCC) attaches proofs of safety
properties to machine code, enabling untrusted code to be verified
by a consumer.  \pcp{} extends this idea to high-level \Python{}
programs with graduated trust and structured certificates.  Unlike
PCC, which targets type safety and memory safety of compiled code,
\pcp{} certifies semantic properties (functional correctness,
relational invariants, resource bounds).

\paragraph{Certified compilation.}
CompCert~(Leroy) and CertiKOS~(Gu et al.)\ produce verified
compilers whose output is guaranteed to preserve the semantics of
the source.  These systems are orthogonal to \pcp{}: they certify
the compiler, not the program.  A verified compiler could be used
to compile \pcp{}-certified \Python{} to certified machine code,
combining both guarantees.

\paragraph{Contract-based verification.}
\CrossHair{} (Buontempo) performs symbolic execution of \Python{}
functions against contracts written as \texttt{assert} statements
or type annotations.  \Nagini{} (Eilers and M\"uller) verifies
\Python{} programs using Viper as a backend.  Both tools produce
binary verdicts without certificates, graduated trust, or
provenance tracking.

\paragraph{Runtime verification.}
\Eiffel{} pioneered Design by Contract, embedding preconditions,
postconditions, and invariants in the language.  iContract and
PyContracts bring similar ideas to Java and \Python{}.  These
systems check contracts at runtime but do not produce certificates
and cannot attest to static properties.

\paragraph{\Dafny{} and \Python{}.}
\Dafny{} can compile verified programs to \Python{}, but the
output is not natural \Python{}: it uses \Dafny{} runtime
libraries, naming conventions, and control flow.  The compiled
code carries no proof information.  \pcp{} verifies natural
\Python{} and ships certificates alongside it.

\paragraph{Dependent types and extraction.}
\Fstar{}, \LEAN{}, and \Coq{} verify programs via dependent types
and extract to host languages.  As discussed in \cref{sec:extraction},
extraction severs the proof-code link and excludes \Python{} as a
target.  \pcp{} avoids extraction entirely.

\paragraph{Semantic sites and descent.}
The \JuGeo{} framework (Papers~01--08 in this series) provides the
mathematical foundations for \pcp{}: semantic sites (Paper~01),
judgment algebra (Paper~02), descent and obstructions (Paper~03),
trust algebra (Paper~04), SMT dispatch (Paper~05), semantic moves
(Paper~06), effect handling (Paper~07), and treaty synthesis
(Paper~08).  This paper builds on all of these to produce
certificate chains.

% =====================================================================
% SECTION 11: CONCLUSION
% =====================================================================
\section{Conclusion}\label{sec:conclusion}

We have presented \emph{Proof-Carrying \Python{}} (\pcp{}), a
methodology in which \JuGeo{} generates certificate chains that
travel with verified \Python{} programs.  Certificate chains
are structured, hash-linked sequences of entries, each recording
a coordinate, proposition, evidence, trust level, hash, and
provenance.  We proved three theoretical results: certificate
soundness (\cref{thm:soundness}), chain compositionality
(\cref{thm:compositionality}), and re-verification completeness
(\cref{thm:reverify}).

The evaluation on \expTotalPrograms{} programs across
\expDomainCount{} domains confirms that \JuGeo{} generates
certificates for all programs with \expObstructionTotal{}
obstructions.  Certificate overhead averages $1.80\times$ the
source size---comparable to extraction overhead but carrying
machine-checkable proof information.  Re-verification completes
in under 16\,$\mu$s, five orders of magnitude faster than
proof-assistant replay.

\pcp{} has three properties that no extraction-based system offers:
\begin{enumerate}
  \item \textbf{Certificates travel with code.}  The proof
    accompanies the program as a JSON artifact, enabling any
    consumer to verify it.
  \item \textbf{Graduated trust.}  Each proposition carries
    a trust level from the \JuGeo{} trust algebra, enabling
    consumers to assess confidence at fine granularity.
  \item \textbf{Zero extraction gap.}  The verified artifact is
    the shipped artifact; no translation, compilation, or
    extraction intervenes between proof and execution.
\end{enumerate}

\paragraph{Future work.}
We plan to extend \pcp{} in three directions.  First, we will
investigate \emph{trust promotion}: using runtime witnesses
($\Truntime$) and SMT solvers ($\Tsolver$) to promote
$\Tcopilot$-level propositions to higher trust, producing
certificates with mixed trust levels.  Second, we will develop
\emph{certificate composition} across package boundaries,
enabling a library's certificate to compose with an application's
certificate.  Third, we will explore \emph{certificate-aware
deployment}: CI/CD pipelines that reject code whose certificate
chain is invalid or whose trust floor falls below a configurable
threshold.

\paragraph{Artifact availability.}
The \JuGeo{} framework, benchmark suite, and all certificate
chains are available at \url{https://github.com/microsoft/jugeo}.

\end{document}
